\chapter{系统信息}
\section{获取CPU信息}
\begin{lstlisting}[style=demostyle,numbers=none]
>>> import psutil
>>> psutil.cpu_count() # CPU逻辑数量
20
>>> psutil.cpu_count(logical=False) # CPU物理核心
12
# 12说明是双核超线程, 20则是20核非超线程
\end{lstlisting}

\section{获取内存信息}
\begin{lstlisting}[style=demostyle,numbers=none]
>>> import psutil
>>> psutil.virtual_memory() # 获取内存信息
svmem(total=33997406208, available=18719956992, percent=44.9, used=15277449216, free=18719956992)
# total: 总内存, available: 可用内存, percent: 使用率, used: 已用内存, free: 空闲内存
>>> psutil.swap_memory() # 获取交换分区信息
sswap(total=2147483648, used=64000000, free=2083483648, percent=3.0, sin=0, sout=0)
# total: 总交换分区, used: 已用交换分区, free: 空闲交换分区, percent: 使用率, sin: 交换分区输入, sout: 交换分区输出
\end{lstlisting}

\section{获取磁盘信息}
\begin{lstlisting}[style=demostyle,numbers=none]
>>> import psutil
>>> psutil.disk_partitions() # 磁盘分区信息
[sdiskpart(device='C:\\', mountpoint='C:\\', fstype='NTFS', opts='rw,fixed'), sdiskpart(device='D:\\', mountpoint='D:\\', fstype='NTFS', opts='rw,fixed'), sdiskpart(device='G:\\', mountpoint='G:\\', fstype='FAT32', opts='rw,fixed')]
# device: 设备名, mountpoint: 挂载点, fstype: 文件系统类型, opts: 挂载选项

>>> psutil.disk_usage('/') # 磁盘使用情况

sdiskusage(total=511243677696, used=457558290432, free=53685387264, percent=89.5)
# total: 总容量, used: 已用容量, free: 空闲容量, percent: 使用率

>>> psutil.disk_io_counters() # 磁盘IO
sdiskio(read_count=7223356, write_count=3159065, read_bytes=272073475072, write_bytes=253278097920, read_time=8198, write_time=4701)
# read_count: 读次数, write_count: 写次数, read_bytes: 读字节数, write_bytes: 写字节数, read_time: 读时间, write_time: 写时间
\end{lstlisting}

\section{获取网络信息}
\begin{lstlisting}[style=demostyle,numbers=none]
>>> import psutil
>>> psutil.net_io_counters() # 获取网络读写字节／包的个数
snetio(bytes_sent=1097595849, bytes_recv=185204061, packets_sent=1197785, packets_recv=412550, errin=0, errout=0, dropin=0, dropout=0)
# bytes_sent: 发送字节数, bytes_recv: 接收字节数, packets_sent: 发送包数, packets_recv: 接收包数, errin: 错误接收数, errout: 错误发送数, dropin: 丢弃接收数, dropout: 丢弃发送数

>>> psutil.net_if_addrs() # 获取网络接口信息
{'以太网': 
    [snicaddr(family=<AddressFamily.AF_LINK: -1>, 
    address='E0-BE-03-96-F1-A1', 
    netmask=None, 
    broadcast=None, 
    ptp=None), 

    snicaddr(family=<AddressFamily.AF_INET: 2>, 
    address='192.168.10.101', 
    netmask='255.255.255.0', 
    broadcast=None, ptp=None), 
    snicaddr(family=<AddressFamily.AF_INET6: 23>, 
    address='2409:8a18:601:5334::49e', 
    netmask=None, 
    broadcast=None, 
    ptp=None), 

    snicaddr(family=<AddressFamily.AF_INET6: 23>, 
    address='2409:8a18:601:5334:c401:e201:52de:444d', 
    netmask=None, 
    broadcast=None, 
    ptp=None), 
    
    snicaddr(family=<AddressFamily.AF_INET6: 23>, 
    address='fdd5:5fdd:d72f::49e', 
    netmask=None, 
    broadcast=None, 
    ptp=None), 
    
    snicaddr(family=<AddressFamily.AF_INET6: 23>, 
    address='fe80::3d12:1c86:bef:a4b', 
    netmask=None, 
    broadcast=None, 
    ptp=None)], 
    
'VMware Network Adapter VMnet1': 
    [snicaddr(family=<AddressFamily.AF_LINK: -1>, 
    address='00-50-56-C0-00-01', 
    netmask=None, 
    broadcast=None, 
    ptp=None), 
    snicaddr(family=<AddressFamily.AF_INET: 2>, 
    address='192.168.159.1', 
    netmask='255.255.255.0', 
    broadcast=None, 
    ptp=None), 
    snicaddr(family=<AddressFamily.AF_INET6: 23>, 
    address='fe80::c89b:fef2:cb75:2d15', 
    netmask=None, 
    broadcast=None, 
    ptp=None)], 
    
'VMware Network Adapter VMnet8': 
    [snicaddr(family=<AddressFamily.AF_LINK: -1>, 
    address='00-50-56-C0-00-08', 
    netmask=None, 
    broadcast=None, 
    ptp=None), 
    snicaddr(family=<AddressFamily.AF_INET: 2>, 
    address='192.168.226.1', 
    netmask='255.255.255.0', 
    broadcast=None, 
    ptp=None), 
    snicaddr(family=<AddressFamily.AF_INET6: 23>, address='fe80::d4c9:92d1:1156:ccd3', netmask=None, broadcast=None, ptp=None)], 

'Loopback Pseudo-Interface 1': 
    [snicaddr(family=<AddressFamily.AF_INET: 2>, address='127.0.0.1', netmask='255.0.0.0', broadcast=None, ptp=None), 
    snicaddr(family=<AddressFamily.AF_INET6: 23>, address='::1', netmask=None, broadcast=None, ptp=None)]
}

>>> psutil.net_if_stats() # 获取网络接口状态
{
    '以太网': snicstats(isup=True, duplex=<NicDuplex.NIC_DUPLEX_FULL: 2>, speed=1000, mtu=1500, flags=''), 
    'VMware Network Adapter VMnet1': snicstats(isup=True, duplex=<NicDuplex.NIC_DUPLEX_FULL: 2>, speed=100, mtu=1500, flags=''), 
    'VMware Network Adapter VMnet8': snicstats(isup=True, duplex=<NicDuplex.NIC_DUPLEX_FULL: 2>, speed=100, mtu=1500, flags=''), 
    'Loopback Pseudo-Interface 1': snicstats(isup=True, duplex=<NicDuplex.NIC_DUPLEX_FULL: 2>, speed=1073, mtu=1500, flags='')
}
\end{lstlisting}

\subsection{网络连接信息}
\begin{lstlisting}[style=demostyle,numbers=none]
>>> import psutil
>>> psutil.net_connections()
[
    sconn(fd=83, family=<AddressFamily.AF_INET6: 30>, type=1, laddr=addr(ip='::127.0.0.1', port=62911), raddr=addr(ip='::127.0.0.1', port=3306), status='ESTABLISHED', pid=3725),
    sconn(fd=84, family=<AddressFamily.AF_INET6: 30>, type=1, laddr=addr(ip='::127.0.0.1', port=62905), raddr=addr(ip='::127.0.0.1', port=3306), status='ESTABLISHED', pid=3725),
    sconn(fd=93, family=<AddressFamily.AF_INET6: 30>, type=1, laddr=addr(ip='::', port=8080), raddr=(), status='LISTEN', pid=3725),
    sconn(fd=103, family=<AddressFamily.AF_INET6: 30>, type=1, laddr=addr(ip='::127.0.0.1', port=62918), raddr=addr(ip='::127.0.0.1', port=3306), status='ESTABLISHED', pid=3725),
    sconn(fd=105, family=<AddressFamily.AF_INET6: 30>, type=1, ..., pid=3725),
    sconn(fd=106, family=<AddressFamily.AF_INET6: 30>, type=1, ..., pid=3725),
    sconn(fd=107, family=<AddressFamily.AF_INET6: 30>, type=1, ..., pid=3725),
    ...
    sconn(fd=27, family=<AddressFamily.AF_INET: 2>, type=2, ..., pid=1)
]
\end{lstlisting}
